\subsection{The thirteen-arrow structural reduction}

\begin{lemma}\label{C:thirteen-three-SCC}
A three-object category with at most thirteen arrows and a
nontrivial strongly connected component cannot have classifying
space of type $\mathbb{CP}^2$.
\end{lemma}
\begin{proof}
An object that is a retract of another can be deleted without
changing the classifying space; the result has at most twelve
arrows and is excluded by Theorem~\ref{C:CP2-lower13}.
Assume there is no such retract.
For a strongly connected three-object category, the finite
categorical Rees description gives finite-group type if all row
numbers or all column numbers are one. Otherwise the two total
numbers are at least four, so its minimal ideal has at least
sixteen arrows, with all three identities outside that ideal.
Both alternatives exclude the proposed type.

The remaining possibility is a strongly connected block on two
objects and one singleton block. The first block has at least six
arrows: each endomorphism monoid has at least two elements and
there is an arrow in each cross direction. A cross arrow between
the blocks yields, by composition inside the first block, at least
one cross arrow for each of its two objects. Thus the first block
has at most ten arrows. A two-object block without a retract that
is not of finite-group type requires at least nine arrows in its
minimal ideal and two identities outside it, hence at least eleven.
Our block is therefore of finite-group type and rationally acyclic.

The singleton endomorphism monoid has order at most five. If its
order is at most four, it too is rationally acyclic. The mixed-nerve
bar decomposition then gives a rational cohomology ring equivalence
from the whole category to a suspension, so every positive-degree
product is zero. If its order is five, every count is forced:
the first block has six arrows and there are exactly two cross
arrows, one for each of its objects. These singleton cross hom-sets
make the category the categorical join of its two blocks, in one
order or the other. Its realization is the join of two connected
spaces, homotopy equivalent to a suspension of their smash product.
Its positive-degree products vanish as well.
\end{proof}

\begin{lemma}\label{C:thirteen-many-objects}
A connected category with thirteen arrows and at least seven
objects has contractible classifying space.
\end{lemma}
\begin{proof}
With $k$ objects connectedness requires at least $k-1$ nonidentity
cross arrows, so there are at least $2k-1$ arrows in all. Thus
$k=7$, there are exactly six cross arrows, and there are no
nonidentity endomorphisms. The underlying unoriented graph must
be a tree, with no parallel or reverse pair: either would leave
too few distinct edges for connectedness. There cannot be a
composable pair of cross arrows, since its composite would give
an additional edge joining the endpoints of a length-two tree
path. The nerve is consequently just that tree, with degeneracies,
and its realization is contractible.
\end{proof}

\begin{remark}\label{C:thirteen-search-scope}
The remaining thirteen-arrow computation for three through six
objects uses \texttt{c\_categories13\_minimal.cpp}. It enumerates
all connected hom-set cardinalities, retains one orbit under
object permutations and opposition, and solves every typed
associativity equation by propagation and exhaustive branching.
At each completed table it keeps the mod-two conditions
$H_1=0$, $\dim H_2=1$, and writes one representative under
identity-fixing permutations within each hom-set.
The latter reduction uses actual relabelings of the entire table;
no partial-table isomorphism heuristic is used.

The only extra deletions are a retract, an initial or terminal
object, the proved singleton-comma deletion, the singleton-join
cases, and Lemma~\ref{C:thirteen-three-SCC} for three objects.
Each remaining representative is checked independently for its
cup square, followed by third and fourth mod-two homology when
the square is nonzero. Unfinished output prefixes are explicitly
marked incomplete and support no complete-census claim. The
executable rejects two-object input; that case uses the separate
triangular and strongly connected arguments.
\end{remark}

\begin{proposition}\label{C:thirteen-four-plus}
A category with thirteen arrows and at least four objects cannot
have classifying space homotopy equivalent to $\mathbb{CP}^2$.
\end{proposition}
\begin{proof}
Disconnected categories cannot have the connected target type.
Lemma~\ref{C:thirteen-many-objects} handles seven or more objects.
For four, five and six objects, the complete algorithm and proved
pruning rules of Remark~\ref{C:thirteen-search-scope} give the
following finished counts. Object hom-cardinality profiles are
already taken modulo relabeling and opposition; the table counts
retain all labels within each such profile.
\[
\begin{array}{c|rrrrr}
\text{objects}&\text{tables}&\text{Hom orbits}
 &\text{classes}&x^2=0&b_3\ne0\\\hline
4&23649952&20692&20123&20103&20\\
5&12902&434&414&414&0\\
6&52&4&3&3&0.
\end{array}
\]
All retained categories have $H_1(-;\mathbb F_2)=0$ and
$\dim H_2(-;\mathbb F_2)=1$. The last two columns exhaust them.
The zero squares have explicit normalized cochain primitives;
the other twenty have nonzero third mod-two homology, computed
using both adjacent boundary ranks. The files
\path{code/outputs/c_categories_13_4_minimal_CP2_cups.json},
\path{code/outputs/c_categories_13_5_minimal_CP2_cups.json}, and
\path{code/outputs/c_categories_13_6_minimal_CP2_cups.json}
are all marked complete, with row counts matching the final
nonempty generator summaries. Their corresponding programs,
input tables and output summaries are supplied. Thus this proof
uses complete enumerations, not a prefix of the still-running
three-object search.
\end{proof}
